\subsubsection{Perception 2D}

For determening the type, and possibly, the distance of objects we use a YOLO26 object classification model. We made our own dataset by recording video using the USB Camera included in the MIRTE Master hardware package. After recording the video we extracted one in every fifteen frames to use for our training dataset. A picture of 4 frames in different environments is shown below

\begin{figure}[!htbp]
\centering
\includegraphics[width=0.7\linewidth]{files/af458f197da824164636251818009f3f.png}
\caption[]{Picture of 4 individual frames}
\label{fig_trainingdataset}
\end{figure}

Training was done manually using Label Studio. Inittially, {\textasciitilde}280 frames were captured in 4 different environments. Due to a large amount of blurring in the images, about 100 pictures were deleted. Of the remaining 180 pictures, 24 pictures were imported into Label Studio and labelled per colour: Green, Red, Purple, White, Light grey, Dark grey, and Black. In total 509 bounding boxes were drawn. The relevant results of the training session are shown below.

\begin{figure}[!htbp]
\centering
\includegraphics[width=0.7\linewidth]{files/4ee2fe1eae7993355b85632bcbc7a4d6.jpeg}
\caption[]{Picture of tests performed during training}
\label{fig_trainingdatasetresults}
\end{figure}

\begin{figure}[!htbp]
\centering
\includegraphics[width=0.7\linewidth]{files/759577a4e723063155d921695fc15039.png}
\caption[]{Plots of trainingresults}
\label{fig_trainingdatasetplots}
\end{figure}

\begin{figure}[!htbp]
\centering
\includegraphics[width=0.7\linewidth]{files/f8ce65a5f1e1e34b97820b1f2a8191e4.png}
\caption[]{Confusion matrix normalized}
\label{fig_trainingdatasetplots}
\end{figure}

\paragraph{Why was YOLO26 used}

The YOLO object-detection family is well established, with YOLO26 being the most recent generation. YOLO26 was chosen because our robot has to detect and classify objects from a live camera on the MIRTE Master platform, which has limited processing power because it uses a OrangePi 3B. Object detection with limited processing power requires a lightweight model with fast inference times while still maintaining accuracy. YOLO26 is suitable for this because it was designed with deployment on low-power devices in mind. YOLO26 changed the post-processing pipeline to NMS-free inference, simplified bounding-box predictions and supports deployment formats commonly used for low-power devices \citep{yolo26keyarchitecturalenhancements}, in our case, the RK3566 NPU of the OrangePi 3B can be used for improving the performance, reducing the total inference time \citep{Yolo11RK3566}. For our model, we chose to use the Nano variant of YOLO26 to minimize inference times. Another benefit of the newer YOLO26 variant over the older variants is that YOLO26 allows for oriented bounding boxes. Unlike standard bouhnding boxes, oriented bounding boxes also estimate the rotation of recognized objects \citep{2026ultralyticsyoloevolutionoverview}. This could be useful for complex grasping tasks where orientation of the gripper is crucial. However, this project only uses standard object detection.